\section{A small cocycle with a prescribed word value}
\label{sec:cocycle}

Let $n\geq 3$ and $\Gamma=\F_n$, with free generators
$x_1,\ldots,x_n$.  We use the left regular action on coefficient
vectors,
\[
  (\lambda(g)h)(b)=h(g^{-1}b),
  \qquad g,b\in\Gamma,
\]
and write $\tau_\Gamma(X)=\langle X\e,\e\rangle$ for the canonical
trace on $L(\Gamma)$.  Given finitely supported real vectors
$h_1,\ldots,h_n$, let $D$ be the cocycle determined by
\[
  D_{x_j}=h_j,
  \qquad D_{gb}=D_g+\lambda(g)D_b.
\]
The construction in this section makes both $h_1$ and
$D_w-\e$ small in $\ell^2(\Gamma)$, with all the other generator
coefficients zero. At time one, Proposition~\ref{prop:flow} converts
these two coefficient bounds into operator-norm bounds for the generator
displacement and the error between the moved word and $CA_w$.
The word $w$ and the support of $h_1$ are allowed to depend on the
required accuracy.
This coefficient rule is the cocycle form of Fox's free differential
calculus~\cite[Section~1, equations~$(1.2)'$, (1.5)--(1.6), and
Section~2, equations~(2.1)--(2.2)]{fox1953}:
successive letters contribute their generator coefficients translated
by the preceding prefix.  We use that rule explicitly below.

\begin{lemma}[Small coefficients with a prescribed word value]
\label{lem:small-cocycle}
Let $n\geq 3$ and $\eta>0$.
There exist a word $w\in\Gamma$ and a finitely supported real
coefficient vector $h$ such that the cocycle determined by
\[
  h_1=h,\qquad h_j=0\quad(2\leq j\leq n)
\]
satisfies
\begin{equation}
\label{eq:small-cocycle}
  \|h\|_{\ell^2}<\eta,
  \qquad
  \|D_w-\e\|_{\ell^2}<\eta.
\end{equation}
\end{lemma}

\subsection{Free prefix words}

For $j\geq 1$, set
\begin{equation}
\label{eq:cocycle-prefixes}
  a=x_1,
  \qquad b_j=x_2^j x_3 x_2^{-j},
  \qquad p_0=e,
  \qquad p_j=(a b_1)\cdots(a b_j).
\end{equation}
For $m\geq 1$ also set
\begin{equation}
\label{eq:cocycle-word-operator}
  w_m=p_m a,
  \qquad T_m=1+\sum_{j=1}^m\lambda(p_j).
\end{equation}

\begin{lemma}[Free prefixes and the cocycle value]
\label{lem:cocycle-prefixes}
For each $m\geq 1$, the elements $p_1,\ldots,p_m$ freely generate a
free subgroup of $\Gamma$.  If $h_1=h$ and $h_j=0$ for $j\geq 2$,
then
\begin{equation}
\label{eq:cocycle-word-value}
  D_{w_m}=T_m h.
\end{equation}
\end{lemma}

\begin{proof}
First consider a nonempty reduced word in $b_1,\ldots,b_m$.
Combining consecutive powers with the same index writes it as
\[
  b_{j_1}^{q_1}\cdots b_{j_l}^{q_l},
  \qquad q_i\in\mathbb Z\setminus\{0\},
  \qquad j_i\ne j_{i+1}.
\]
Its expansion in $x_2,x_3$ is
\[
  x_2^{j_1}x_3^{q_1}
  x_2^{j_2-j_1}x_3^{q_2}\cdots
  x_2^{j_l-j_{l-1}}x_3^{q_l}x_2^{-j_l}.
\]
Every displayed $x_3$ power is nonzero, and each intervening $x_2$
power is nonzero.  This word cannot reduce to the identity.
Thus the $b_j$ are a free basis of the subgroup they generate.
The free product decomposition
\[
  \Gamma=\langle a\rangle *\langle x_2,\ldots,x_n\rangle
\]
then shows that $a,b_1,\ldots,b_m$ freely generate their subgroup:
an alternating product of nonidentity elements from
$\langle a\rangle$ and $\langle b_1,\ldots,b_m\rangle$ is also
reduced in the displayed free product.

In the abstract free group on $a,b_1,\ldots,b_m$, the substitution
fixing $a$ and sending each $b_j$ to $a b_j$ has inverse given by
$b_j\mapsto a^{-1}b_j$, again fixing $a$.
Consequently
\[
  u_j=a b_j,\qquad 1\leq j\leq m,
\]
form a free family.  On the abstract free group with free basis
$u_1,\ldots,u_m$, define two endomorphisms by
\[
  \varphi(u_j)=u_1\cdots u_j,
  \qquad
  \psi(u_1)=u_1,
  \qquad
  \psi(u_j)=u_{j-1}^{-1}u_j\quad (j\geq 2).
\]
For every $j$, cancellation gives
\begin{align*}
  \psi\varphi(u_j)
    &=u_1(u_1^{-1}u_2)\cdots(u_{j-1}^{-1}u_j)=u_j,\\
  \varphi\psi(u_j)
    &=(u_1\cdots u_{j-1})^{-1}(u_1\cdots u_j)=u_j
       \quad (j\geq 2),
\end{align*}
and the second identity also holds for $j=1$.
Thus $\varphi$ and $\psi$ are inverse automorphisms.
Since $p_j=u_1\cdots u_j$, this proves the assertion about the
$p_j$.

For the cocycle assertion, the inverse rule
$D_{x_j^{-1}}=-\lambda(x_j^{-1})D_{x_j}$ shows that $D$ vanishes
on the subgroup generated by $x_2,\ldots,x_n$.
In particular, $D_{b_j}=0$.
The positive occurrences of $a$ in $w_m$ have respective prefixes
$p_0,p_1,\ldots,p_m$, and there are no other occurrences of $a$
or $a^{-1}$.
Applying the cocycle identity successively therefore gives
\[
  D_{w_m}=\sum_{j=0}^m\lambda(p_j)h=T_mh.
\]
\end{proof}

\subsection{The spectral measures of the prefix sums}

Put
\begin{equation}
\label{eq:cocycle-positive-operator}
  Q_m=T_mT_m^*,
  \qquad
  \mu_m(B)=\tau_\Gamma\bigl(1_B(Q_m/m)\bigr)
\end{equation}
for Borel sets $B\subseteq[0,\infty)$.
These are probability measures by the spectral theorem and the
normalization of $\tau_\Gamma$.

We seek a small vector $h$ with $T_mh$ close to $\e$. The spectral
limit below has no atom at zero. It will let us choose a small fixed
cutoff in the normalized spectrum and then make $m$ large enough that
the discarded projection has small trace. On the remaining spectrum,
the inverse of $Q_m$ has norm at most a fixed constant times $m^{-1}$.
The inverse bound controls the coefficient norm, while the discarded
trace controls the approximation error.
The Catalan count below is the moment calculation associated with
circular systems in Voiculescu's free-probability framework;
see~\cite{voiculescu1985} and~\cite[Definition~1.9 and Remark~1.11]{voiculescu1990}.
We give the word count and the normalization of the limiting density
explicitly.

\begin{lemma}[Spectral limit of the prefix sums]
\label{lem:cocycle-spectrum}
The measures $\mu_m$ converge weakly to the probability measure
\begin{equation}
\label{eq:cocycle-limit-density}
  d\nu(x)=\frac{1}{2\pi}\sqrt{\frac{4-x}{x}}
           \,1_{(0,4)}(x)\,dx.
\end{equation}
In particular, $\nu(\{0\})=0$.
\end{lemma}

\begin{proof}
Write $V_j=\lambda(p_j)$ and $V_0=1$.
Lemma~\ref{lem:cocycle-prefixes} implies that a word in
$V_1,\ldots,V_m$ and their inverses has trace one if its word in
these free symbols reduces to the identity, and trace zero
otherwise.
This trace rule applies inside the given regular representation.
Indeed, for $H=\langle p_1,\ldots,p_m\rangle$, the spaces
$\ell^2(Ht)$, as $Ht$ ranges over the right cosets of $H$, are
reducing copies of the left regular representation of $H$.
The vector $\e$ lies in the copy $\ell^2(H)$.
Its vector state therefore agrees with the canonical trace of
$H$ on polynomials and on the spectral projections under
consideration.

Fix an integer $r\geq 1$.  Expanding $Q_m^r$ gives
\begin{equation}
\label{eq:cocycle-moment-expansion}
  \tau_\Gamma(Q_m^r)
  =\sum_{i_1,\ldots,i_{2r}\in\{0,\ldots,m\}}
    \tau_\Gamma\bigl(
      V_{i_1}V_{i_2}^*\cdots V_{i_{2r-1}}V_{i_{2r}}^*
    \bigr).
\end{equation}
Each summand is zero or one.  In a summand equal to one, every
nonzero index used must occur at least twice, since a generator
occurring only once cannot disappear under free reduction.
There are therefore at most $r$ distinct nonzero indices.

If exactly $r$ distinct nonzero indices occur, all $2r$ positions
are nonidentity and each index occurs exactly twice.
The two occurrences have opposite exponents, as follows already
by abelianizing the identity word.
Pair the two positions with the same index.  This matching is
unique for the chosen word.
A sequence of adjacent inverse cancellations produces a
noncrossing matching: a pair deleted after other pairs can only
enclose pairs already deleted, and two deleted pairs cannot
cross.
Conversely, a noncrossing matching with opposite exponents at
its endpoints reduces to the identity by deleting an innermost
pair and proceeding by induction.
Every noncrossing perfect matching of $\{1,\ldots,2r\}$ has
opposite parity at the two endpoints of each pair.  To see this,
the positions strictly inside that pair must be matched among
themselves, so their number is even.
The signs in~\eqref{eq:cocycle-moment-expansion} alternate;
therefore every noncrossing perfect matching has the required
opposite signs.

The number of noncrossing perfect matchings is
\[
  \Cat_r=\frac{1}{r+1}\binom{2r}{r}.
\]
For completeness, splitting a matching at the partner of its
first position gives the recurrence
$c_0=1$ and $c_r=\sum_{j=0}^{r-1}c_jc_{r-1-j}$.
Its formal generating series satisfies $C(z)=1+zC(z)^2$,
so the solution with constant term one is
$(1-\sqrt{1-4z})/(2z)$; the binomial expansion gives the displayed
coefficient $\Cat_r$.
For each matching, labeling its $r$ pairs with distinct indices
from $\{1,\ldots,m\}$ gives
\[
  (m)_r=m(m-1)\cdots(m-r+1)
\]
choices, with the value zero when $m<r$.
Ordering the pairs by their first endpoints makes this count
unambiguous, and the uniqueness of the equal-index matching
prevents double counting.
Thus the contribution from exactly $r$ distinct indices is
$\Cat_r(m)_r$.

Every remaining contributing term has at most $r-1$ distinct
nonzero indices.  This includes every term with an identity
summand: such a term has at most $2r-1$ nonidentity positions,
while each index it uses occurs at least twice.
For fixed $r$, there are only finitely many patterns specifying
the identity positions and the partition of the other positions
into equal-index classes.  A pattern with $k\leq r-1$ classes
has at most $m^k$ labelings.  Summing over those patterns gives
an $O_r(m^{r-1})$ bound, with its constant depending only on $r$.
We have proved
\begin{equation}
\label{eq:cocycle-catalan-moments}
  \tau_\Gamma(Q_m^r)
    =\Cat_r(m)_r+O_r(m^{r-1}),
  \qquad
  \int x^r\,d\mu_m(x)\longrightarrow\Cat_r.
\end{equation}

We next justify passage from moments to weak convergence.
For $r=1$, direct expansion gives
$\tau_\Gamma(Q_m)=m+1$, so the first moments of $\mu_m$ are
bounded by two.  Markov's inequality gives tightness on
$[0,\infty)$.
For each fixed integer $r\geq 1$, the moment convergence in
\eqref{eq:cocycle-catalan-moments}, applied also at $r+1$, gives
\[
  \sup_m\int x^{r+1}\,d\mu_m(x)<\infty,
  \qquad
  \sup_m\int_{(R,\infty)}x^r\,d\mu_m(x)
    \leq \frac{1}{R}\sup_m\int x^{r+1}\,d\mu_m(x)
    \longrightarrow 0
\]
as $R\to\infty$.
Consequently, if a subsequence converges weakly to a probability
measure $\mu$, its moments pass to the limit.
More explicitly, apply weak convergence to the bounded continuous
functions $\min\{x^r,R^r\}$ and use the displayed tail bound;
the corresponding bound for $\mu$ follows by first applying
weak convergence to bounded truncations of $x^{r+1}$.
Letting $R\to\infty$ proves
\begin{equation}
\label{eq:cocycle-subsequence-moments}
  \int x^r\,d\mu(x)=\Cat_r\qquad(r\geq 0).
\end{equation}

The density in~\eqref{eq:cocycle-limit-density} realizes these
moments, including the zeroth moment. Write
$B(a,b)=\int_0^1t^{a-1}(1-t)^{b-1}\,dt$ for the Euler beta integral.
The substitution
$x=4t$ and the beta integral give, for every integer $r\geq 0$,
\begin{align*}
  \frac{1}{2\pi}\int_0^4 x^{r-1/2}(4-x)^{1/2}\,dx
    &=\frac{4^{r+1}}{2\pi}
      B\left(r+\frac12,\frac32\right)\\
    &=\frac{(2r)!}{r!(r+1)!}
      =\Cat_r.
\end{align*}
The middle equality follows from
$B(\tfrac12,\tfrac32)=\pi/2$ and the recurrence
$B(a+1,b)=aB(a,b)/(a+b)$.
In particular the density has total mass one.

These moments determine a probability measure on
$[0,\infty)$ uniquely.
Indeed, $\Cat_r\leq 4^r$.  If a measure satisfying
\eqref{eq:cocycle-subsequence-moments} had mass $c>0$ on
$[4+\epsilon,\infty)$ for some $\epsilon>0$, then
\[
  4^r\geq\Cat_r\geq c(4+\epsilon)^r
\]
for every $r$, a contradiction for large $r$.
Its support is therefore contained in $[0,4]$.
On this compact interval, uniform polynomial approximation of
continuous functions shows that equality of all moments implies
equality of the measures.
Every subsequential weak limit is consequently $\nu$.
Tightness now gives weak convergence of the full sequence
$\mu_m$ to $\nu$.
Finally, $\nu$ has the displayed integrable density and thus has
no atom at zero.
\end{proof}

\subsection{Cutoff inversion and finite coefficients}



\begin{proof}[Proof of Lemma~\ref{lem:small-cocycle}]
Use the operators $T_m,Q_m$ and the measures $\mu_m$ above.
Since $\nu(\{0\})=0$, choose a fixed $d>0$ so small that
\[
  \nu([0,d])<\frac{\eta^2}{16}.
\]
The interval $[0,d]$ is closed, so weak convergence and the
Portmanteau inequality give
\[
  \limsup_{m\to\infty}\mu_m([0,d])
     \leq\nu([0,d])<\frac{\eta^2}{16}.
\]
We may therefore choose $m$ sufficiently large that both
\begin{equation}
\label{eq:cocycle-cutoff-choices}
  \mu_m([0,d])<\frac{\eta^2}{8},
  \qquad
  \frac{1}{dm}<\frac{\eta^2}{16}
\end{equation}
hold.  Keep this $d$ and this $m$ fixed for the remainder of the
construction.

Set
\[
  E_m=1_{(dm,\infty)}(Q_m),
  \qquad R_m=f_m(Q_m),
  \qquad
  f_m(x)=
  \begin{cases}
    0,&0\leq x\leq dm,\\
    x^{-1},&x>dm.
  \end{cases}
\]
Thus $R_m$ is a bounded positive operator, with
$\norm{R_m}\leq (dm)^{-1}$.
The spectral theorem gives
\begin{equation}
\label{eq:cocycle-cutoff-identities}
  Q_m R_m=E_m,
  \qquad R_m Q_m R_m=R_m,
  \qquad
  \tau_\Gamma(1-E_m)=\mu_m([0,d]).
\end{equation}
In particular, these operations allow a kernel or arbitrarily
small nonzero spectrum of $Q_m$.
Define the coefficient vector
\[
  h^{(0)}=T_m^*R_m\e\in\ell^2(\Gamma).
\]
Using $Q_m=T_mT_m^*$ in the indicated order gives
\begin{align}
\label{eq:cocycle-cutoff-target}
  T_mh^{(0)}&=E_m\e,\\
\label{eq:cocycle-cutoff-error}
  \|T_mh^{(0)}-\e\|_{\ell^2}^2
    &=\tau_\Gamma(1-E_m)<\frac{\eta^2}{8},\\
\label{eq:cocycle-cutoff-size}
  \|h^{(0)}\|_{\ell^2}^2
    &=\langle R_mQ_mR_m\e,\e\rangle
      =\tau_\Gamma(R_m)
      \leq\frac{1}{dm}<\frac{\eta^2}{16}.
\end{align}

The vector $h^{(0)}$ has real coefficients.
To check this directly, let $J$ denote coefficientwise complex
conjugation on $\ell^2(\Gamma)$.
Each left shift is real in the standard basis, so $J$ commutes
with $T_m,T_m^*$ and $Q_m$, and $J\e=\e$.
Commutation with real polynomials in $Q_m$ extends by uniform
approximation to real continuous functions of $Q_m$.
The continuous functions
\[
  x\longmapsto \min\{1,\max\{0,l(x-dm)\}\},
  \qquad l=1,2,\ldots,
\]
converge pointwise, with a uniform bound, to
$1_{(dm,\infty)}(x)$ on the spectrum of $Q_m$.
Their functional calculi converge strongly to $E_m$, hence
$J$ also commutes with $E_m$.
Finally,
\[
  R_m=E_m g_m(Q_m),
  \qquad g_m(x)=\frac{1}{\max\{dm,x\}},
\]
where $g_m$ is real and continuous on $[0,\infty)$.
Thus $J$ commutes with $R_m$ as well, proving
$Jh^{(0)}=h^{(0)}$.

Real finitely supported coefficient vectors are dense in the
real subspace of $\ell^2(\Gamma)$.
Choose such a vector $h$ with
\[
  \|h-h^{(0)}\|_{\ell^2}
       <\frac{\eta}{4(m+1)}.
\]
The elementary bound $\norm{T_m}\leq m+1$ and
\eqref{eq:cocycle-cutoff-error}--\eqref{eq:cocycle-cutoff-size}
yield
\begin{align*}
  \|h\|_{\ell^2}
    &<\frac{\eta}{4}+\frac{\eta}{4(m+1)}<\eta,\\
  \|T_mh-\e\|_{\ell^2}
    &<\frac{\eta}{\sqrt{8}}+\frac{\eta}{4}<\eta.
\end{align*}
For this finitely supported $h$, form the cocycle with
$h_1=h$ and all other generator values zero, and take $w=w_m$.
Equation~\eqref{eq:cocycle-word-value} gives $D_w=T_mh$, proving
\eqref{eq:small-cocycle}.

The choices have been made in the order $d$, then $m$, then the
finite coefficient approximation.  The last approximation uses
the norm of the already fixed operator $T_m$; the lemma requires
no bound on the support size of $h$ uniform in $\eta$.
\end{proof}
